% Conclusions. Built by the future pass from the instructions below. Do not
% process now.

[Write a focused, two to three paragraph conclusion that restates the contribution
and the forward path without introducing new data; reference Results figures and
tables rather than repeating numbers. Open by restating the headline finding from
\texttt{papers/VVUQ-02/execution/README.md}: an autonomous Claude Code Opus 4.7 1M
Max run executed every entry point, passed all 172 tests, drove the full ten-gate
ACCEPT, BLOCK, and ESCALATE surface, cleared all 32 sweep iterations, processed the
1790-line tournament and the 1000-row sensor stream, and reproduced all three large
artifacts from seed 20260525, all under one pull request. State plainly that this
operationalises the thesis: for a single autonomous surgical humanoid the assurance
process, not code generation, is the substantial and decision-bearing work, and
holding VVUQ to a higher standard than the code is what the record shows.]

[Close by arguing that binding every gate to published external standards already
used in real life is what turns an inexpensive autonomous run into defensible
evidence, that this checks real assurance steps off so the field does not have to
re-derive them for each deliverable, and that adapting the approach across oncology
clinical trials promises faster, less expensive, and more patient-beneficial
deliverables than conventional verification. End respectfully but opportunistically
on the regulatory opportunity and on US leadership in surgical-humanoid VVUQ tied to
external standards, and reaffirm that the H2-Surgical 1.0 is a hypothetical 2030
platform whose every candidate behavior must clear the ten gates and a recorded
human reviewer before any non-simulated use. Cite the prior framework
(\texttt{kawchak2026vvuq01}) and the credibility standards (\texttt{asmevv40},
\texttt{nasastd7009a}) once each.]
